\documentclass[10pt,a4paper]{amsart}
\usepackage{amsmath,amssymb,amsthm,graphicx}
% Retained minimal preamble for the Stacks Project chapter more-algebra.tex (More on Algebra).
% The chapter file pinned-original-source.tex carries only \input{preamble} on its first line, so no
% preamble text lives inside the chapter itself.  The Stacks Project uses one single shared preamble for
% every chapter; the snapshot stages it identically as sources/stacks-chow/preamble.tex and
% sources/stacks-etale-cohomology/preamble.tex.  The block below is copied mechanically, line for line,
% from sources/stacks-chow/preamble.tex lines 14-18 (the Xy-pic setup) and 154-217 (the native theorem
% declarations together with the mathematical macros).  Nothing was added, removed or redefined; the
% wrapper class and the amsmath/amssymb/amsthm/graphicx base packages are supplied by the assembler.
% Chapter: more-algebra.tex -- More on Algebra.
% For commutative diagrams we use Xy-pic
\usepackage[all]{xy}

% We use 2cell for 2-commutative diagrams.
\xyoption{2cell}

% Theorem environments.
%
\theoremstyle{plain}
\newtheorem{theorem}[subsection]{Theorem}
\newtheorem{proposition}[subsection]{Proposition}
\newtheorem{lemma}[subsection]{Lemma}

\theoremstyle{definition}
\newtheorem{definition}[subsection]{Definition}
\newtheorem{example}[subsection]{Example}
\newtheorem{exercise}[subsection]{Exercise}
\newtheorem{situation}[subsection]{Situation}

\theoremstyle{remark}
\newtheorem{remark}[subsection]{Remark}
\newtheorem{remarks}[subsection]{Remarks}

\numberwithin{equation}{subsection}

% Macros
%
\def\lim{\mathop{\mathrm{lim}}\nolimits}
\def\colim{\mathop{\mathrm{colim}}\nolimits}
\def\Spec{\mathop{\mathrm{Spec}}}
\def\Hom{\mathop{\mathrm{Hom}}\nolimits}
\def\Ext{\mathop{\mathrm{Ext}}\nolimits}
\def\SheafHom{\mathop{\mathcal{H}\!\mathit{om}}\nolimits}
\def\SheafExt{\mathop{\mathcal{E}\!\mathit{xt}}\nolimits}
\def\Sch{\mathit{Sch}}
\def\Mor{\mathop{\mathrm{Mor}}\nolimits}
\def\Ob{\mathop{\mathrm{Ob}}\nolimits}
\def\Sh{\mathop{\mathit{Sh}}\nolimits}
\def\NL{\mathop{N\!L}\nolimits}
\def\CH{\mathop{\mathrm{CH}}\nolimits}
\def\proetale{{pro\text{-}\acute{e}tale}}
\def\etale{{\acute{e}tale}}
\def\QCoh{\mathit{QCoh}}
\def\Ker{\mathop{\mathrm{Ker}}}
\def\Im{\mathop{\mathrm{Im}}}
\def\Coker{\mathop{\mathrm{Coker}}}
\def\Coim{\mathop{\mathrm{Coim}}}

% Boxtimes
%
\DeclareMathSymbol{\boxtimes}{\mathbin}{AMSa}{"02}

%
% Macros for moduli stacks/spaces
%
\def\QCohstack{\mathcal{QC}\!\mathit{oh}}
\def\Cohstack{\mathcal{C}\!\mathit{oh}}
\def\Spacesstack{\mathcal{S}\!\mathit{paces}}
\def\Quotfunctor{\mathrm{Quot}}
\def\Hilbfunctor{\mathrm{Hilb}}
\def\Curvesstack{\mathcal{C}\!\mathit{urves}}
\def\Polarizedstack{\mathcal{P}\!\mathit{olarized}}
\def\Complexesstack{\mathcal{C}\!\mathit{omplexes}}
% \Pic is the operator that assigns to X its picard group, usage \Pic(X)
% \Picardstack_{X/B} denotes the Picard stack of X over B
% \Picardfunctor_{X/B} denotes the Picard functor of X over B
\def\Pic{\mathop{\mathrm{Pic}}\nolimits}
\def\Picardstack{\mathcal{P}\!\mathit{ic}}
\def\Picardfunctor{\mathrm{Pic}}
\def\Deformationcategory{\mathcal{D}\!\mathit{ef}}
\usepackage{cleveref}
\crefname{originallemma}{original lemma}{original lemmas}
\Crefname{originallemma}{Original lemma}{Original lemmas}
\makeatletter
\newlabel{lemma-hom-from-Af}{{091P}{source}{}{sourcepointer-1}{}}
\newlabel{lemma-hom-from-Af@cref}{{[originallemma][0][]091P}{[1][0][]source}{}{}{}}
\makeatother
\crefname{originallemma}{original lemma}{original lemmas}
\Crefname{originallemma}{Original lemma}{Original lemmas}
\makeatletter
\newlabel{algebra-lemma-when-surjective-to-completion}{{090S}{source}{}{sourcepointer-2}{}}
\newlabel{algebra-lemma-when-surjective-to-completion@cref}{{[originallemma][0][]090S}{[1][0][]source}{}{}{}}
\makeatother
\title{I-adic completeness of a module is equivalent to the vanishing of its torsion functors and surjectivity onto its completion}
\author{The Stacks Project Authors}
\date{stacks-project:more-algebra.tex:local-snapshot-20261012; GNU FDL 1.2 or later}
\begin{document}
\maketitle
\setcounter{subsection}{1}
\noindent\textit{Editorial note.} All delivered mathematical text is the author's own, copied from the pinned Stacks Project chapter file more-algebra.tex; no statement or proof step was rewritten, reordered or removed. The counted statement is the labeled lemma environment at source lines 25962-25970 (canonical tag 091R) and the counted proof is the proof environment at source lines 25972-26015, which immediately follows it and is its own proof.\par\medskip
\noindent\textit{Editorial note.} The proof refers to lemma-hom-from-Af inside this chapter and to algebra-lemma-when-surjective-to-completion of another chapter; both are navigation pointers with their roles recorded in the external reference reasons. The infinite correction series of part (1) and the universal extension of part (2) are inside the counted span.\par\medskip
\noindent\textit{Editorial note.} The proof is divided by the author into the two numbered parts with medskip headings, and both are delivered in full.\par\medskip
\noindent\textit{Editorial note.} The chapter file pinned-original-source.tex carries only an input of the shared preamble on its first line, so no preamble text lives inside the chapter. The retained file sources/stacks-more-algebra/retained-preamble.tex is a mechanical, line-for-line copy of the project's single shared preamble as staged in the snapshot (Xy-pic setup plus the native theorem declarations and mathematical macros). Nothing was added, removed or redefined, and the excerpt serves only as preamble, never as a span coordinate.\par\medskip
\noindent\textit{Editorial note.} The source URL is the source-level Stacks Project URL, as required when more than one card is produced from one source; the per-claim tag is carried by claim\_key and by this note. The counted statement and proof cite no bibliography entry, so this card needs no bibliography.\par\medskip
\noindent\textit{Allowed standard prerequisites (editorial).}\begin{itemize}
\item For an element f of a ring, the torsion functor T(-, f) is computed by the Hom and Ext groups from the localization A\_f, so its vanishing is equivalent to the vanishing of Hom\_A(A\_f, M) and Ext\textasciicircum{}1\_A(A\_f, M). invoked at proof lines 25974-25975 and again at proof lines 26005-26006. Turns the vanishing condition on the torsion functor into two explicit vanishing statements about maps out of A\_f, which are then proved by hand.
\item An I-adically complete module is the inverse limit of its quotients M over I\textasciicircum{}n, and Hom from A\_f into a module annihilated by a power of f vanishes. used at proof lines 25976-25977. Gives the vanishing of Hom\_A(A\_f, M) for a complete module.
\item The snake lemma identifies the quotients M over f\textasciicircum{}n M and E over f\textasciicircum{}n E for an extension of M by A\_f, and a short exact sequence of modules in which the outer terms are killed by a power of f splits. used at proof lines 26012-26013 and at proof lines 26005-26007. Produces the congruence expressing a given element x as a finite combination of the x\_n modulo f\textasciicircum{}n M, which is the surjectivity assertion of part (2).
\end{itemize}\medskip
\ifdefined\hypertarget\hypertarget{sourcepointer-1}{}\fi
\noindent\textit{Source reference.} Original lemma 091P: sources/stacks-more-algebra/pinned-original-source.tex, source lines 25850--25850. The reference identifies the vanishing of T(M, f) with the vanishing of Hom and Ext from A\_f, used at proof lines 25974-25975 and 26005-26006; the two explicit constructions that verify those vanishing conditions are inside the counted span.\par\medskip
\ifdefined\hypertarget\hypertarget{sourcepointer-2}{}\fi
\noindent\textit{Source reference.} Original lemma 090S: sources/stacks-chow/tags.txt, source lines 11420--11420. The reference is the reduction step that lets the finitely generated case be replaced by the case of one element, used at proof lines 25993-25995; the extension built afterwards and the snake lemma argument are inside the counted span.\par\medskip
\begin{lemma}
\label{lemma-complete-derived-complete}
Let $A$ be a ring. Let $I \subset A$ be an ideal. Let $M$ be an $A$-module.
\begin{enumerate}
\item If $M$ is $I$-adically complete, then $T(M, f) = 0$ for all $f \in I$.
\item Conversely, if $T(M, f) = 0$ for all $f \in I$ and $I$ is finitely
generated, then $M \to \lim M/I^nM$ is surjective.
\end{enumerate}
\end{lemma}

\begin{proof}
Proof of (1). Assume $M$ is $I$-adically complete.
By Lemma \ref{lemma-hom-from-Af} it suffices to prove
$\Ext^1_A(A_f, M) = 0$ and $\Hom_A(A_f, M) = 0$.
Since $M = \lim M/I^nM$ and since $\Hom_A(A_f, M/I^nM) = 0$
it follows that $\Hom_A(A_f, M) = 0$. Suppose we have an extension
$$
0 \to M \to E \to A_f \to 0
$$
For $n \geq 0$ pick $e_n \in E$ mapping to $1/f^n$.
Set $\delta_n = fe_{n + 1} - e_n \in M$ for $n \geq 0$.
Replace $e_n$ by
$$
e'_n = e_n + \delta_n + f\delta_{n + 1} + f^2 \delta_{n + 2} + \ldots
$$
The infinite sum exists as $M$ is complete with respect to $I$ and $f \in I$.
A simple calculation shows that $fe'_{n + 1} = e'_n$. Thus we get a splitting
of the extension by mapping $1/f^n$ to $e'_n$.

\medskip\noindent
Proof of (2). Assume that $I = (f_1, \ldots, f_r)$ and that $T(M, f_i) = 0$
for $i = 1, \ldots, r$. By
Algebra, Lemma \ref{algebra-lemma-when-surjective-to-completion}
we may assume $I = (f)$ and $T(M, f) = 0$. Let $x_n \in M$ for $n \geq 0$.
Consider the extension
$$
0 \to M \to E \to A_f \to 0
$$
given by
$$
E = M \oplus \bigoplus Ae_n\Big/\langle x_n - fe_{n + 1} + e_n\rangle
$$
mapping $e_n$ to $1/f^n$ in $A_f$ (see above).
By assumption and Lemma \ref{lemma-hom-from-Af}
this extension is split, hence we obtain an element
$x + e_0$ which generates a copy of $A_f$ in $E$.
Then
$$
x + e_0 = x - x_0 + fe_1 = x - x_0 - f x_1 + f^2 e_2 = \ldots
$$
Since $M/f^nM = E/f^nE$ by the snake lemma, we see that
$x = x_0 + fx_1 + \ldots + f^{n - 1}x_{n - 1}$ modulo $f^nM$.
In other words, the map $M \to \lim M/f^nM$ is surjective as desired.
\end{proof}

\end{document}
